% `tex.web` §1010: how much room a `\vsplit`-able insertion class is given.
%
%     if count(n)<=0 then w:=max_dimen
%     else begin w:=page_goal-page_total-page_depth;
%       if count(n)<>1000 then w:=x_over_n(w,count(n))*1000;
%       end;
%     if w>dimen(n)-height(r) then w:=dimen(n)-height(r);
%
% The scaling divides *first* and multiplies after — the same shape as §1009's
% `h:=x_over_n(height(p),1000)*count(n)` — and the order is not
% interchangeable. `xn_over_d(w,1000,count)` is exact where this deliberately
% is not: it drops the remainder before scaling it up, so the answer is always
% a multiple of 1000sp and always a little short of the exact quotient.
%
% We used the exact form. Measured against real TeX, four independent inputs,
% all off by the rounding and nothing else:
%
%     w        count   TeX (x_over_n*1000)   ours (xn_over_d)
%     90pt     500     179.99268             180.0
%     93pt     500     185.98938             186.0
%     100pt    500     199.99695             200.0
%     100pt    250     399.9939              400.0
%
% This was trip.tex's first page-trace divergence. The original investigation
% measured the gap at 120sp and — assuming the arithmetic was exact — concluded
% it must be in `page_depth`, the one
% quantity in the expression that nothing prints. The arithmetic was the gap.
% The note is deleted with this change; it was blocked on reading §1010, and
% varying the inputs against real TeX answered it without needing to.
\catcode`\{=1 \catcode`\}=2
\tracingpages=1 \tracingonline=0
\output={\shipout\box255}
\vsize=100pt \maxdepth=2pt
\dimen200=10000pt \count200=500 \skip200=3pt

% An insertion far too tall for the page, so §1010 has to decide how much of it
% fits. Two firings, so two splits with different `page_total` behind them.
\hbox{}\vskip 7pt
\insert200{\hbox{}\vskip 400pt\hbox{}\vskip 400pt}
\hbox{}\vskip 300pt
\end
